\documentclass[11pt,a4paper]{article}

\usepackage[utf8]{inputenc}
\usepackage[margin=1in]{geometry}
\usepackage{amsmath,amssymb}
\usepackage{graphicx}
\usepackage{booktabs}
\usepackage{tabularx}
\usepackage{caption}
\usepackage{float}
\usepackage{hyperref}
\usepackage{fancyhdr}
\usepackage{titlesec}
\usepackage{xcolor}
\usepackage{listings}
\usepackage{tcolorbox}
\usepackage{enumitem}

% Hyperref setup
\hypersetup{
    colorlinks=true,
    linkcolor=blue!70!black,
    citecolor=blue!70!black,
    urlcolor=blue!70!black
}

% Header and Footer setup
\pagestyle{fancy}
\fancyhf{}
\fancyhead[L]{\small INT315 -- Cluster Computing Project Report}
\fancyhead[R]{\small Apache Spark MLlib Student Dropout Prediction}
\fancyfoot[C]{\thepage}
\renewcommand{\headrulewidth}{0.4pt}
\renewcommand{\footrulewidth}{0.4pt}

% Code Listing Styling
\definecolor{codegreen}{rgb}{0,0.6,0}
\definecolor{codegray}{rgb}{0.5,0.5,0.5}
\definecolor{codepurple}{rgb}{0.58,0,0.82}
\definecolor{backcolour}{rgb}{0.97,0.98,0.99}

\lstset{
    backgroundcolor=\color{backcolour},
    commentstyle=\color{codegreen},
    keywordstyle=\color{blue}\bfseries,
    numberstyle=\tiny\color{codegray},
    stringstyle=\color{codepurple},
    basicstyle=\ttfamily\footnotesize,
    breakatwhitespace=false,
    breaklines=true,
    captionpos=b,
    keepspaces=true,
    numbers=left,
    numbersep=5pt,
    showspaces=false,
    showstringspaces=false,
    showtabs=false,
    tabsize=4,
    frame=single,
    rulecolor=\color{gray!30}
}

\begin{document}

% ==============================================================================
% 1. TITLE & COVER PAGE (Requirement 1)
% ==============================================================================
\begin{titlepage}
    \centering
    \vspace*{0.5cm}
    
    {\Huge \bfseries INT315: CLUSTER COMPUTING \par}
    \vspace{0.4cm}
    {\Large \textsc{Department of Computer Science \& Engineering} \par}
    {\large \textsc{Academic Session: 2026--2027} \par}
    
    \vspace{1.2cm}
    \rule{\linewidth}{1.5pt} \\[0.4cm]
    {\huge \bfseries Distributed Early Warning System for Student Dropout Risk Prediction using Apache Spark MLlib on Cluster Infrastructure \par}
    \rule{\linewidth}{1.5pt} \\[1.5cm]
    
    \begin{minipage}[t]{0.48\textwidth}
        \textbf{\large Submitted By:} \\[0.2cm]
        \textbf{Name:} {[Student Name]} \\
        \textbf{Registration No:} {[Registration Number]} \\
        \textbf{Roll No:} {[Roll Number]} \\
        \textbf{Section / Group:} {[Section]} \\
        \textbf{Degree:} B.Tech Computer Science and Engineering
    \end{minipage}
    \hfill
    \begin{minipage}[t]{0.48\textwidth}
        \begin{flushright}
            \textbf{\large Submitted To:} \\[0.2cm]
            \textbf{Course Instructor / Evaluator:} \\[0.1cm]
            {[Faculty / Instructor Name]} \\
            \textit{Course Code: INT315} \\
            \textit{School of Computer Science \& Engineering} \\[0.1cm]
            {[University / Institute Name]}
        \end{flushright}
    \end{minipage}
    
    \vfill
    
    \begin{tcolorbox}[colback=blue!5!white,colframe=blue!75!black,title=\centering \bfseries Academic Declaration]
        \small This project report is submitted in partial fulfillment of the requirements for the course \textbf{INT315 -- Cluster Computing}. It encompasses distributed data engineering, resilient MLlib pipeline orchestration, cluster-level classification benchmarks, and real-time interactive decision support.
    \end{tcolorbox}
    
    \vspace{0.5cm}
    {\large \today \par}
\end{titlepage}

\tableofcontents
\newpage

% ==============================================================================
% 2. PROBLEM STATEMENT (Requirement 2)
% ==============================================================================
\section{Problem Statement}
\label{sec:problem_statement}

In higher education institutions worldwide, student attrition and academic dropout present severe academic, institutional, and economic challenges. Withdrawing prior to degree completion leads to wasted instructional capacity, financial forfeiture for students and funding bodies, and diminished institutional reputation. Identifying students at imminent risk of attrition in early semesters is crucial, as targeted institutional interventions---such as individualized academic tutoring, counseling, and financial bursary outreach---are significantly more effective when applied before irreversible academic failure occurs.

However, modern universities collect vast, heterogeneous datasets spanning demographic backgrounds, secondary school qualifications, enrollment modes, semester-by-semester curricular evaluations, fee payment histories, and macroeconomic context. Processing, cleaning, transforming, and modeling these high-dimensional educational datasets poses substantial computational bottlenecks when executed on traditional single-node machines:
\begin{enumerate}[label=(\alph*)]
    \item \textbf{Memory Saturation \& Ingestion Bottlenecks}: High-cardinality multi-semester educational records quickly exhaust single-node RAM during multi-stage feature encoding, cross-validation, and recursive tree construction.
    \item \textbf{Combinatoric Tree Search Overhead}: Training ensemble tree classifiers (such as Random Forest) across continuous academic metrics (grades, attendance, evaluations) requires massive parallel sorting and candidate split evaluations across feature subsets, leading to severe latency on standard CPUs.
    \item \textbf{Class Imbalance \& High-Stakes Classification}: Academic dropouts represent a critical minority class ($\sim 32\%$). Conventional single-node models often bias predictions toward the majority graduating class, failing to provide the high recall (sensitivity) required for an actionable early warning system.
\end{enumerate}

Therefore, there is a compelling need for a \textbf{distributed cluster computing architecture} leveraging \textbf{Apache Spark 4.2.0 (PySpark MLlib)} to perform parallelized feature vector assembly, distributed classification modeling, and scalable real-time risk stratification.

% ==============================================================================
% 3. OBJECTIVE OF THE PROJECT (Requirement 3)
% ==============================================================================
\section{Objective of the Project}
\label{sec:objective}

The primary objective of this project is to architect, implement, evaluate, and deploy an end-to-end distributed machine learning system for student dropout risk prediction on cluster computing infrastructure. The specific technical objectives are:

\begin{enumerate}
    \item \textbf{Distributed Data Ingestion \& Cleansing (Task 97)}: Ingest high-dimensional student records into an Apache Spark distributed DataFrame using resilient partitioning; execute distributed null handling and map multi-class targets into a binary classification schema ($Y=1.0$ for Dropout, $Y=0.0$ for Retained).
    \item \textbf{Distributed Feature Vector Assembly (Task 98)}: Construct an Apache Spark MLlib feature engineering pipeline utilizing \texttt{StringIndexer} and \texttt{OneHotEncoder} across categorical variables, coupled with a distributed \texttt{VectorAssembler} to yield a dense, scaled 43-dimensional feature vector.
    \item \textbf{Dual Distributed Classification Modeling (Task 99)}: Formulate, train, and serialize two distinct cluster-compatible MLlib pipelines:
    \begin{itemize}
        \item \textbf{Logistic Regression}: Distributed gradient descent optimization with elastic-net regularization.
        \item \textbf{Random Forest Classifier}: Distributed parallel ensemble tree construction utilizing distributed histogram aggregation and bagging.
    \end{itemize}
    \item \textbf{Rigorous Cluster Evaluation \& Benchmarking (Task 100)}: Evaluate trained models on an unseen 20\% holdout test partition using distributed evaluators for \textbf{ROC-AUC}, \textbf{F1-Score}, \textbf{Precision}, \textbf{Recall}, and specialized \textbf{Dropout Recall (Sensitivity)}.
    \item \textbf{Continuous Risk Scoring \& Stratification (Task 101)}: Extract continuous dropout probability vectors ($P(\text{Dropout})$) and stratify the holdout cohort into three actionable early warning tiers: \textbf{High Risk} ($P \ge 0.60$), \textbf{Moderate Risk} ($0.35 \le P < 0.60$), and \textbf{Low Risk} ($P < 0.35$).
    \item \textbf{Interactive Decision Support Dashboard Deployment}: Deploy a multi-tab \textbf{Streamlit + Plotly} early warning system directly satisfying evaluator requirements:
    \begin{itemize}
        \item View 1: Student Risk Categories \& Probability Distributions.
        \item View 2: Actual vs. Predicted Outcomes with interactive Confusion Matrix Heatmap.
        \item View 3: Dropout Rates segmented across all 17 decoded Academic Programs.
        \item View 4: Real-time PySpark model serving for live student risk inference.
        \item View 5: Advisor outreach roster with one-click CSV export for targeted interventions.
    \end{itemize}
\end{enumerate}

% ==============================================================================
% 4. DATASET DESCRIPTION (Requirement 4)
% ==============================================================================
\section{Dataset Description}
\label{sec:dataset_description}

The system utilizes the authoritative \textbf{Predict Students' Dropout and Academic Success} dataset. The dataset comprises \textbf{4,424 student records} characterized across \textbf{35 distinct attributes} captured at the time of university enrollment and during the first and second academic semesters.

\subsection{Feature Taxonomy}
The attributes are logically categorized into four functional domains:
\begin{enumerate}
    \item \textbf{Curricular Academic Performance (Primary Predictors)}:
    \begin{itemize}
        \item \texttt{Curricular units 1st sem (evaluations, approved, grade)}: First-semester academic workload, course pass rates, and grade averages (0--20 scale).
        \item \texttt{Curricular units 2nd sem (evaluations, approved, grade)}: Second-semester examination attendance, approved units, and GPA.
        \item \texttt{Previous qualification (grade)}: Entrance qualification grades.
        \item \texttt{Admission grade}: Comprehensive score obtained during university matriculation.
    \end{itemize}
    \item \textbf{Demographic Indicators}:
    \begin{itemize}
        \item \texttt{Age at enrollment}: Student chronological age (range 17--60, median 20).
        \item \texttt{Gender}: Binary indicator ($1=\text{Male}, 0=\text{Female}$).
        \item \texttt{Marital status}: Categorical code ($1=\text{Single}, 2=\text{Married}, \dots$).
        \item \texttt{Displaced student}: Whether the student relocated geographically to attend university.
    \end{itemize}
    \item \textbf{Socioeconomic \& Institutional Factors}:
    \begin{itemize}
        \item \texttt{Course}: Academic program of study (17 distinct degree programs).
        \item \texttt{Daytime/evening attendance}: Attendance mode ($1=\text{Daytime}, 0=\text{Evening}$).
        \item \texttt{Scholarship holder}: Binary indicator of state financial aid.
        \item \texttt{Tuition fees up to date}: Financial fee payment status ($1=\text{Current}, 0=\text{Arrears}$).
        \item \texttt{Debtor}: Outstanding unpaid institutional balances.
        \item \texttt{Mother's/Father's qualification and occupation}: Parental educational attainment and socioeconomic status.
    \end{itemize}
    \item \textbf{Macroeconomic Indicators}:
    \begin{itemize}
        \item \texttt{Unemployment rate}, \texttt{Inflation rate}, \texttt{GDP}: Regional economic conditions at the time of enrollment.
    \end{itemize}
\end{enumerate}

\subsection{Target Variable Formulation}
The raw dataset contains a ternary target variable \texttt{Target} categorized into:
\begin{itemize}
    \item \textbf{Dropout}: 1,421 students (32.12\%) --- Students who discontinued their degree studies.
    \item \textbf{Graduate}: 2,209 students (49.93\%) --- Students who successfully completed their degree.
    \item \textbf{Enrolled}: 794 students (17.95\%) --- Students actively pursuing their coursework.
\end{itemize}

For optimal institutional utility and standard probabilistic ROC-AUC evaluation, the target is cast into a \textbf{Binary Classification Problem}:
\begin{equation}
    Y = \begin{cases}
        1.0 & \text{if } \texttt{Target} = \text{``Dropout'' (At-Risk Class)} \\
        0.0 & \text{if } \texttt{Target} \in \{\text{``Graduate''}, \text{``Enrolled''}\} \text{ (Retained Class)}
    \end{cases}
\end{equation}

% ==============================================================================
% 5. DATASET SOURCE (Requirement 5)
% ==============================================================================
\section{Dataset Source}
\label{sec:dataset_source}

The dataset used in this cluster computing project is derived from the established academic study conducted by the Polytechnic Institute of Portalegre, Portugal:

\begin{itemize}
    \item \textbf{Original Academic Publication}: \\
    Realinho, V., Machado, J., Baptista, L., \& Martins, M. V. (2021). \textit{Predicting Student Dropout and Academic Success}. Data, 6(11), 110. \\
    \textbf{DOI}: \url{https://doi.org/10.3390/data6110110}
    
    \item \textbf{Official UCI Machine Learning Repository}: \\
    UC Irvine Machine Learning Repository: \textit{Predict Students' Dropout and Academic Success}. \\
    \textbf{Direct Dataset URL}: \url{https://archive.ics.uci.edu/dataset/697/predict+students+dropout+and+academic+success}
    
    \item \textbf{Kaggle Dataset Mirror}: \\
    Kaggle: \textit{Predict Student Dropout and Academic Success}. \\
    \textbf{Kaggle URL}: \url{https://www.kaggle.com/datasets/thedevastator/higher-education-predictors-of-student-retention}
\end{itemize}

The raw comma-separated values file (\texttt{dataset.csv}) contains 4,424 records and is ingested via Spark's distributed file reader using localized URI paths (\texttt{file:///}).

% ==============================================================================
% 6. TOOLS \& TECHNOLOGIES USED (Requirement 6)
% ==============================================================================
\section{Tools \& Technologies Used}
\label{sec:tools_technologies}

To fulfill the requirements of **INT315 -- Cluster Computing**, the project architecture relies strictly on distributed data processing and machine learning frameworks:

\begin{table}[H]
    \centering
    \caption{Distributed Cluster Computing \& Software Stack}
    \label{tab:tools_stack}
    \small
    \begin{tabularx}{\linewidth}{lX}
        \toprule
        \textbf{Component} & \textbf{Role in Cluster Computing Architecture} \\
        \midrule
        \textbf{Apache Spark 4.2.0} & Distributed general-purpose cluster computing engine providing fault-tolerant in-memory processing, DAG scheduling, and distributed memory abstraction across worker nodes. \\
        \addlinespace
        \textbf{PySpark MLlib} & Native Spark distributed machine learning library. Implements pipeline abstractions (\texttt{Pipeline}, \texttt{PipelineModel}), distributed feature engineering (\texttt{StringIndexer}, \texttt{OneHotEncoder}, \texttt{VectorAssembler}), and parallel model trainers. \\
        \addlinespace
        \textbf{Spark Catalyst Optimizer} & Relational query optimization engine that automatically translates high-level DataFrame transformations into optimized physical execution plans (logical plan $\to$ physical plan $\to$ Java bytecode). \\
        \addlinespace
        \textbf{Project Tungsten} & Spark engine optimization layer utilizing off-heap memory management and whole-stage code generation to eliminate JVM object overhead and maximize CPU cache efficiency. \\
        \addlinespace
        \textbf{Py4J Gateway} & High-throughput IPC bridge enabling dynamic Python interaction with the underlying Spark JVM runtime without serialization penalty. \\
        \addlinespace
        \textbf{OpenJDK 21.0.12.1} & Modern 64-Bit Java Virtual Machine powering Spark Driver and Executor processes. \\
        \addlinespace
        \textbf{Streamlit (v1.66+)} & Modern reactive web application framework serving the serialized PySpark \texttt{PipelineModel} live for clinical/institutional early warning. \\
        \addlinespace
        \textbf{Plotly Express} & GPU-accelerated interactive visualization library for confusion matrices, donut charts, and grouped program attrition rates. \\
        \addlinespace
        \textbf{PyTest (v9.2+)} & Automated test framework validating distributed transformations, schema types, and analytical accuracy via 14 unit and integration tests. \\
        \bottomrule
    \end{tabularx}
\end{table}

\subsection{Cluster Computing Execution Architecture}
The computation follows the distributed master-worker cluster paradigm:
\begin{itemize}
    \item \textbf{Spark Driver}: Hosts the \texttt{SparkSession}, orchestrates job execution, constructs the Directed Acyclic Graph (DAG), and serializes tasks for distributed workers.
    \item \textbf{Cluster Executors}: In-memory worker threads that ingest DataFrame partitions in parallel, execute map/filter transformations, compute partial gradient descent sums, and aggregate distributed decision tree histograms.
    \item \textbf{Lazy Evaluation}: Transformation operations (\texttt{withColumn}, \texttt{filter}) build an execution lineage without execution until an action (\texttt{count}, \texttt{collect}, \texttt{save}) triggers DAG compilation.
\end{itemize}

% ==============================================================================
% 7. IMPLEMENTATION / METHODOLOGY (Requirement 7)
% ==============================================================================
\section{Implementation / Methodology}
\label{sec:methodology}

The project is implemented across five structured tasks adhering to the distributed PySpark lifecycle.

\subsection{Task 97: Distributed Data Ingestion \& Null Cleansing}
The ingestion layer (\texttt{src/data\_loader.py}) reads \texttt{dataset.csv} into a partitioned Spark DataFrame. 
\begin{itemize}
    \item \textbf{Filesystem URI Safety}: Configures \texttt{spark.hadoop.fs.defaultFS = "file:///"} to enforce local filesystem paths, preventing connection failures caused by idle HDFS NameNodes (port 9000).
    \item \textbf{Binary Target Labeling}: Converts \texttt{Target} strings to numeric double labels using Spark SQL expressions:
    \begin{equation}
        \texttt{label} = \text{when}(\texttt{col("Target")} == \text{``Dropout''}, 1.0).\text{otherwise}(0.0)
    \end{equation}
    \item \textbf{Defensive Null Imputation}: Categorical columns are cast to \texttt{StringType} and nulls are defensively imputed with \texttt{"Unknown"} using \texttt{df.na.fill()}.
\end{itemize}

\subsection{Task 98: Categorical Encoding \& Vector Assembly}
Distributed feature transformation (\texttt{src/feature\_pipeline.py}) converts heterogeneous columns into an MLlib-compatible feature vector:
\begin{enumerate}
    \item \textbf{StringIndexer}: Maps distinct string categories across 8 categorical columns (\texttt{Course}, \texttt{Daytime/evening attendance}, \texttt{Marital status}, \texttt{Gender}, \texttt{Scholarship holder}, \texttt{Tuition fees up to date}, \texttt{Debtor}, \texttt{Displaced}) into integer indices $[0, K-1]$.
    \item \textbf{OneHotEncoder}: Transforms indexed categories into sparse binary indicator vectors, preventing artificial ordinal ranking assumptions.
    \item \textbf{VectorAssembler}: Combines the one-hot vectors with 10 continuous numeric features (evaluations, approved units, GPA grades across 1st and 2nd semesters, enrollment age, unemployment, inflation, GDP) into a unified \textbf{43-dimensional Spark \texttt{features} vector}.
\end{enumerate}

\subsection{Task 99: Distributed Classification Modeling}
Two distinct machine learning algorithms are trained via Spark MLlib pipelines (\texttt{src/model\_trainer.py}):

\subsubsection{1. Distributed Logistic Regression}
Minimizes the regularized binary cross-entropy loss across distributed partitions:
\begin{equation}
    L(w) = -\frac{1}{N} \sum_{i=1}^N \left[ y_i \log(\sigma(w^T x_i)) + (1 - y_i) \log(1 - \sigma(w^T x_i)) \right] + \lambda \left( \alpha \|w\|_1 + \frac{1-\alpha}{2} \|w\|_2^2 \right)
\end{equation}
where $\sigma(z) = \frac{1}{1 + e^{-z}}$, $\lambda = 0.01$, and $\alpha = 0.5$ (elastic-net mixing parameter). Spark executors compute partition-level loss gradients and aggregate them via tree-reduce operations.

\subsubsection{2. Distributed Random Forest Classifier}
Builds an ensemble of $B=80$ de-correlated decision trees with maximum tree depth $D=8$. 
At each split node, a random subset of features is sampled. For continuous features, worker nodes construct distributed equi-depth histograms to evaluate candidate split thresholds in parallel, maximizing the Information Gain ($\Delta I$):
\begin{equation}
    \Delta I = H(D) - \left( \frac{|D_L|}{|D|} H(D_L) + \frac{|D_R|}{|D|} H(D_R) \right)
\end{equation}
where $H(D)$ represents the Gini impurity:
\begin{equation}
    H(D) = 1 - \sum_{k=0}^1 p_k^2
\end{equation}
Both trained models are serialized to disk (\texttt{models/lr\_dropout\_model} and \texttt{models/rf\_dropout\_model}) for low-latency batch and online serving.

\subsection{Task 100: Model Evaluation \& Benchmarking}
The holdout test partition (20\%, $N=843$ students) is scored and evaluated using Spark's \texttt{MulticlassClassificationEvaluator} and \texttt{BinaryClassificationEvaluator}:
\begin{itemize}
    \item \textbf{ROC-AUC}: Measures the area under the Receiver Operating Characteristic curve, assessing discrimination ability across all thresholds.
    \item \textbf{F1-Score (Weighted)}: Harmonic mean of precision and recall.
    \item \textbf{Precision \& Recall}: Measures positive predictive value and total coverage.
    \item \textbf{Dropout Recall (Sensitivity)}: Specialized metric evaluating recall strictly on the positive class ($Y=1.0$):
    \begin{equation}
        \text{Dropout Recall} = \frac{\text{True Positives (TP)}}{\text{True Positives (TP)} + \text{False Negatives (FN)}}
    \end{equation}
\end{itemize}

\subsection{Task 101: Test Scoring \& 3-Tier Risk Stratification}
Continuous probability vectors are extracted via Spark's \texttt{vector\_to\_array} function. Each test student is classified into an early-warning risk tier:
\begin{equation}
    \text{Risk Tier} = \begin{cases}
        \text{\textbf{High Risk (Immediate Intervention)}} & \text{if } P(\text{Dropout}) \ge 0.60 \\
        \text{\textbf{Moderate Risk (Active Monitoring)}} & \text{if } 0.35 \le P(\text{Dropout}) < 0.60 \\
        \text{\textbf{Low Risk (On Track)}} & \text{if } P(\text{Dropout}) < 0.35
    \end{cases}
\end{equation}
The scored cohort is exported to \texttt{output/test\_dropout\_predictions.csv} and evaluation metrics to \texttt{output/metrics.json}.

% ==============================================================================
% 8. COMPLETE SOURCE CODE (Requirement 8)
% ==============================================================================
\section{Complete Source Code}
\label{sec:source_code}

The complete, production-grade source code for all six core PySpark cluster modules is provided below:

\subsection{Module 1: SparkSession Factory (\texttt{src/spark\_utils.py})}
\begin{lstlisting}[language=Python, caption={SparkSession configuration with local filesystem enforcement}]
import os
import sys
import warnings

# Suppress harmless pandas / arrow version warnings
warnings.filterwarnings("ignore", category=UserWarning)
warnings.filterwarnings("ignore", category=FutureWarning)

# Ensure Spark and Py4J libraries from /opt/spark are on Python's path
if "/opt/spark/python" not in sys.path:
    os.environ["SPARK_HOME"] = "/opt/spark"
    sys.path.insert(0, "/opt/spark/python/lib/py4j-0.10.9.9-src.zip")
    sys.path.insert(0, "/opt/spark/python")

from pyspark.sql import SparkSession


def get_spark_session(app_name: str = "Student_Dropout_Pipeline") -> SparkSession:
    """
    Initializes and returns a properly configured local SparkSession.
    Configures Spark to use the local filesystem, avoiding unneeded HDFS connection attempts.
    """
    spark = (
        SparkSession.builder
        .appName(app_name)
        .master("local[*]")
        .config("spark.driver.memory", "4g")
        .config("spark.sql.shuffle.partitions", "8")
        .config("spark.hadoop.fs.defaultFS", "file:///")
        .config("spark.sql.execution.arrow.pyspark.enabled", "false")
        .getOrCreate()
    )
    # Set log level to WARN to reduce console clutter
    spark.sparkContext.setLogLevel("WARN")
    return spark

\end{lstlisting}

\subsection{Module 2: Ingestion \& Cleansing (\texttt{src/data\_loader.py})}
\begin{lstlisting}[language=Python, caption={Distributed data ingestion, binary target mapping, and null imputation}]
import os
from typing import List
from pyspark.sql import SparkSession, DataFrame
from pyspark.sql.functions import col, when

# Categorical columns to cast to StringType for StringIndexer
CATEGORICAL_COLS = [
    "Course",
    "Daytime/evening attendance",
    "Marital status",
    "Gender",
    "Scholarship holder",
    "Tuition fees up to date",
    "Debtor",
    "Displaced"
]

TARGET_COL = "label"


def load_student_data(spark: SparkSession, file_path: str) -> DataFrame:
    """
    Task 97: Loads the student dataset into a Spark DataFrame
    using an explicit absolute file:// path.
    """
    abs_path = os.path.abspath(file_path)
    file_uri = f"file://{abs_path}" if not abs_path.startswith("file://") else abs_path

    if not os.path.exists(abs_path):
        raise FileNotFoundError(f"Student dataset not found at: {abs_path}")

    df = spark.read.csv(file_uri, header=True, inferSchema=True)
    return df


def preprocess_student_data(df: DataFrame) -> DataFrame:
    """
    Task 97: Converts Target to binary label (1.0 = Dropout, 0.0 = Retained),
    casts categorical columns to StringType, and handles missing/null values.
    """
    # 1. Binary target mapping
    if "Target" in df.columns:
        df_transformed = df.withColumn(
            "label",
            when(col("Target") == "Dropout", 1.0).otherwise(0.0)
        )
    else:
        df_transformed = df

    # 2. Cast categorical integer codes to strings for StringIndexer
    for c in CATEGORICAL_COLS:
        if c in df_transformed.columns:
            df_transformed = df_transformed.withColumn(c, col(c).cast("string"))

    # 3. Defensive null handling
    fill_string_dict = {c: "Unknown" for c in CATEGORICAL_COLS if c in df_transformed.columns}
    if fill_string_dict:
        df_clean = df_transformed.na.fill(fill_string_dict)
    else:
        df_clean = df_transformed

    # Filter out null targets if present
    if "label" in df_clean.columns:
        df_clean = df_clean.filter(col("label").isNotNull())

    return df_clean


def load_and_preprocess(spark: SparkSession, file_path: str) -> DataFrame:
    """Convenience helper combining data loading and preprocessing."""
    df_raw = load_student_data(spark, file_path)
    return preprocess_student_data(df_raw)

\end{lstlisting}

\subsection{Module 3: Feature Pipeline (\texttt{src/feature\_pipeline.py})}
\begin{lstlisting}[language=Python, caption={StringIndexer, OneHotEncoder, and VectorAssembler pipeline stages}]
from typing import List, Tuple, Any
from pyspark.ml.feature import StringIndexer, OneHotEncoder, VectorAssembler

from src.data_loader import CATEGORICAL_COLS, TARGET_COL

NUMERIC_COLS = [
    "Curricular units 1st sem (evaluations)",
    "Curricular units 1st sem (approved)",
    "Curricular units 1st sem (grade)",
    "Curricular units 2nd sem (evaluations)",
    "Curricular units 2nd sem (approved)",
    "Curricular units 2nd sem (grade)",
    "Age at enrollment",
    "Unemployment rate",
    "Inflation rate",
    "GDP"
]


def build_feature_stages() -> Tuple[List[Any], str]:
    """
    Task 98: Constructs the feature engineering pipeline stages:
    1. StringIndexer: Encodes categorical features (course, study mode, demographics).
    2. OneHotEncoder: Transforms indexed categories into sparse one-hot vectors.
    3. VectorAssembler: Combines numeric metrics (evaluations, approved units, grades)
       and encoded vectors into a single 'features' vector.

    Returns:
        stages: List of feature pipeline stages.
        features_col: Name of the assembled features vector column ('features').
    """
    indexed_cols = [f"{c}_idx" for c in CATEGORICAL_COLS]
    encoded_cols = [f"{c}_vec" for c in CATEGORICAL_COLS]

    # StringIndexer stage
    indexer = StringIndexer(
        inputCols=CATEGORICAL_COLS,
        outputCols=indexed_cols,
        handleInvalid="keep"
    )

    # OneHotEncoder stage
    encoder = OneHotEncoder(
        inputCols=indexed_cols,
        outputCols=encoded_cols
    )

    # VectorAssembler stage
    assembler_inputs = NUMERIC_COLS + encoded_cols
    assembler = VectorAssembler(
        inputCols=assembler_inputs,
        outputCol="features",
        handleInvalid="skip"
    )

    stages = [indexer, encoder, assembler]
    return stages, "features"

\end{lstlisting}

\subsection{Module 4: Model Trainer (\texttt{src/model\_trainer.py})}
\begin{lstlisting}[language=Python, caption={PySpark MLlib Logistic Regression and Random Forest model pipelines}]
import os
from typing import Dict
from pyspark.ml import Pipeline, PipelineModel
from pyspark.ml.classification import LogisticRegression, RandomForestClassifier
from pyspark.sql import DataFrame

from src.feature_pipeline import build_feature_stages, TARGET_COL


def build_logistic_regression_pipeline(
    max_iter: int = 30,
    reg_param: float = 0.01,
    elastic_net_param: float = 0.5
) -> Pipeline:
    """Task 99: Builds a Spark ML Pipeline with Logistic Regression."""
    feature_stages, features_col = build_feature_stages()
    lr = LogisticRegression(
        featuresCol=features_col,
        labelCol=TARGET_COL,
        maxIter=max_iter,
        regParam=reg_param,
        elasticNetParam=elastic_net_param
    )
    return Pipeline(stages=feature_stages + [lr])


def build_random_forest_pipeline(
    num_trees: int = 80,
    max_depth: int = 8,
    seed: int = 42
) -> Pipeline:
    """Task 99: Builds a Spark ML Pipeline with Random Forest Classifier."""
    feature_stages, features_col = build_feature_stages()
    rf = RandomForestClassifier(
        featuresCol=features_col,
        labelCol=TARGET_COL,
        numTrees=num_trees,
        maxDepth=max_depth,
        seed=seed
    )
    return Pipeline(stages=feature_stages + [rf])


def train_models(
    train_df: DataFrame,
    save_dir: str = "models"
) -> Dict[str, PipelineModel]:
    """
    Task 99: Trains both Logistic Regression and Random Forest models
    on training data and serializes them to disk.
    """
    os.makedirs(save_dir, exist_ok=True)

    print("Training Logistic Regression Pipeline...")
    pipeline_lr = build_logistic_regression_pipeline()
    model_lr = pipeline_lr.fit(train_df)
    lr_path = os.path.join(save_dir, "lr_dropout_model")
    model_lr.write().overwrite().save(f"file://{os.path.abspath(lr_path)}")
    print(f"Saved Logistic Regression model to: {lr_path}")

    print("Training Random Forest Classifier Pipeline...")
    pipeline_rf = build_random_forest_pipeline()
    model_rf = pipeline_rf.fit(train_df)
    rf_path = os.path.join(save_dir, "rf_dropout_model")
    model_rf.write().overwrite().save(f"file://{os.path.abspath(rf_path)}")
    print(f"Saved Random Forest model to: {rf_path}")

    return {
        "logistic_regression": model_lr,
        "random_forest": model_rf
    }


def load_model(model_path: str) -> PipelineModel:
    """Loads a serialized PySpark PipelineModel from disk."""
    abs_path = os.path.abspath(model_path)
    return PipelineModel.load(f"file://{abs_path}")

\end{lstlisting}

\subsection{Module 5: Distributed Evaluator (\texttt{src/evaluator.py})}
\begin{lstlisting}[language=Python, caption={Classification evaluation metrics and JSON metrics exporter}]
import json
import os
from typing import Dict, Any
from pyspark.ml.evaluation import MulticlassClassificationEvaluator, BinaryClassificationEvaluator
from pyspark.sql import DataFrame

from src.data_loader import TARGET_COL


def evaluate_classification_predictions(
    predictions_df: DataFrame,
    label_col: str = TARGET_COL,
    prediction_col: str = "prediction",
    raw_pred_col: str = "rawPrediction"
) -> Dict[str, float]:
    """
    Task 100: Evaluates classification performance using:
    - Precision
    - Recall
    - F1-Score
    - ROC-AUC
    - Dropout-specific Recall (Sensitivity on class 1.0)
    """
    eval_f1 = MulticlassClassificationEvaluator(
        labelCol=label_col,
        predictionCol=prediction_col,
        metricName="f1"
    )
    eval_precision = MulticlassClassificationEvaluator(
        labelCol=label_col,
        predictionCol=prediction_col,
        metricName="weightedPrecision"
    )
    eval_recall = MulticlassClassificationEvaluator(
        labelCol=label_col,
        predictionCol=prediction_col,
        metricName="weightedRecall"
    )
    eval_auc = BinaryClassificationEvaluator(
        labelCol=label_col,
        rawPredictionCol=raw_pred_col,
        metricName="areaUnderROC"
    )

    f1 = eval_f1.evaluate(predictions_df)
    precision = eval_precision.evaluate(predictions_df)
    recall = eval_recall.evaluate(predictions_df)
    auc = eval_auc.evaluate(predictions_df)

    # Dropout-specific recall (Sensitivity on at-risk class 1.0)
    eval_dropout_recall = MulticlassClassificationEvaluator(
        labelCol=label_col,
        predictionCol=prediction_col,
        metricName="recallByLabel",
        metricLabel=1.0
    )
    dropout_recall = eval_dropout_recall.evaluate(predictions_df)

    return {
        "roc_auc": round(auc, 4),
        "f1_score": round(f1, 4),
        "precision": round(precision, 4),
        "recall": round(recall, 4),
        "dropout_recall": round(dropout_recall, 4)
    }


def compare_and_save_metrics(
    eval_results: Dict[str, Any],
    save_path: str = "output/metrics.json"
) -> None:
    """Saves model evaluation comparison to a JSON file."""
    os.makedirs(os.path.dirname(save_path), exist_ok=True)
    with open(save_path, "w") as f:
        json.dump(eval_results, f, indent=4)
    print(f"Metrics saved to: {save_path}")

\end{lstlisting}

\subsection{Module 6: Pipeline Runner Orchestrator (\texttt{src/pipeline\_runner.py})}
\begin{lstlisting}[language=Python, caption={End-to-end pipeline execution, test scoring, and 3-tier risk stratification}]
import os
import sys

from src.spark_utils import get_spark_session
from pyspark.sql.functions import col, when
from pyspark.ml.functions import vector_to_array

from src.data_loader import load_and_preprocess
from src.model_trainer import train_models
from src.evaluator import evaluate_classification_predictions, compare_and_save_metrics


def run_pipeline(
    data_path: str = "src/data/dataset.csv",
    output_dir: str = "output",
    models_dir: str = "models"
) -> None:
    """
    Executes Tasks 97 through 101 end-to-end:
    - Task 97: Load data & handle missing values
    - Task 98: Feature vector assembly
    - Task 99: Train Logistic Regression & Random Forest Classifier
    - Task 100: Evaluate using Precision, Recall, F1-Score, and ROC-AUC
    - Task 101: Score holdout test set & categorize dropout risk tiers
    """
    print("=" * 70)
    print("[ACADEMIC] Starting Student Dropout Risk PySpark Pipeline (Tasks 97-101)")
    print("=" * 70)

    # Resolve data path (supports src/data/dataset.csv or data/dataset.csv)
    if not os.path.exists(data_path) and os.path.exists("data/dataset.csv"):
        data_path = "data/dataset.csv"

    # 1. Initialize SparkSession
    spark = get_spark_session("Student_Dropout_Runner")

    # 2. Task 97: Data Loading & Preprocessing
    print(f"\n[Task 97] Loading student records from: {data_path}")
    df_clean = load_and_preprocess(spark, data_path)
    total_count = df_clean.count()
    dropout_count = df_clean.filter(col("label") == 1.0).count()
    print(f"[Task 97] Cleaned records: {total_count:,} (At-Risk Dropouts: {dropout_count:,} / {dropout_count/total_count*100:.1f}%)")

    # 3. Train-Test Split (80 / 20)
    print("\nSplitting dataset into 80% Train and 20% Holdout Test...")
    train_df, test_df = df_clean.randomSplit([0.8, 0.2], seed=42)
    train_count = train_df.count()
    test_count = test_df.count()
    print(f"Train samples: {train_count:,} | Test samples: {test_count:,}")

    # 4. Tasks 98 & 99: Train Classification Models
    print("\n[Tasks 98 & 99] Training classification models via Spark MLlib...")
    models = train_models(train_df, save_dir=models_dir)

    # 5. Task 100: Comprehensive Model Evaluation
    print("\n[Task 100] Evaluating models on holdout test set...")
    preds_lr = models["logistic_regression"].transform(test_df)
    preds_rf = models["random_forest"].transform(test_df)

    metrics_lr = evaluate_classification_predictions(preds_lr)
    metrics_rf = evaluate_classification_predictions(preds_rf)

    print("\n" + "=" * 65)
    print("[TARGET] Model Performance Metrics (Holdout Test Set)")
    print("=" * 65)
    print(f"{'Metric':<38} {'Logistic Reg':<14} {'Random Forest':<14}")
    print("-" * 65)
    print(f"{'ROC-AUC (Discrimination Power)':<38} {metrics_lr['roc_auc']:<14.4f} {metrics_rf['roc_auc']:<14.4f}")
    print(f"{'F1-Score (Balanced Accuracy)':<38} {metrics_lr['f1_score']:<14.4f} {metrics_rf['f1_score']:<14.4f}")
    print(f"{'Precision (Positive Predictive)':<38} {metrics_lr['precision']:<14.4f} {metrics_rf['precision']:<14.4f}")
    print(f"{'Recall (Overall Sensitivity)':<38} {metrics_lr['recall']:<14.4f} {metrics_rf['recall']:<14.4f}")
    print(f"{'Dropout Recall (At-Risk Catch)':<38} {metrics_lr['dropout_recall']:<14.4f} {metrics_rf['dropout_recall']:<14.4f}")
    print("=" * 65)

    all_metrics = {
        "logistic_regression": metrics_lr,
        "random_forest": metrics_rf,
        "total_students": total_count,
        "train_samples": train_count,
        "test_samples": test_count
    }
    metrics_path = os.path.join(output_dir, "metrics.json")
    compare_and_save_metrics(all_metrics, save_path=metrics_path)

    # 6. Task 101: Test Set Risk Scoring & Tier Categorization
    print("\n[Task 101] Scoring holdout students and stratifying into early warning risk tiers...")
    scored_rf = preds_rf.withColumn("dropout_probability", vector_to_array(col("probability"))[1])

    # 3-Tier Risk Assignment
    scored_rf = scored_rf.withColumn(
        "risk_level",
        when(col("dropout_probability") >= 0.60, "High Risk")
        .when(col("dropout_probability") >= 0.35, "Moderate Risk")
        .otherwise("Low Risk")
    )

    export_cols = [
        "Course",
        "Daytime/evening attendance",
        "Scholarship holder",
        "Tuition fees up to date",
        "Debtor",
        "Age at enrollment",
        "Curricular units 1st sem (evaluations)",
        "Curricular units 1st sem (approved)",
        "Curricular units 1st sem (grade)",
        "Curricular units 2nd sem (evaluations)",
        "Curricular units 2nd sem (approved)",
        "Curricular units 2nd sem (grade)",
        "dropout_probability",
        "risk_level",
        "prediction",
        "label"
    ]

    os.makedirs(output_dir, exist_ok=True)
    csv_out_path = os.path.join(output_dir, "test_dropout_predictions.csv")
    pdf_export = scored_rf.select(export_cols).toPandas()
    pdf_export.to_csv(csv_out_path, index=False)
    print(f"[Task 101] Successfully exported {len(pdf_export):,} scored student predictions to: {csv_out_path}")

    # Risk Distribution Breakdown
    risk_summary = pdf_export["risk_level"].value_counts()
    print("\n[ROSTER] Early Warning Cohort Distribution:")
    for risk_tier, count in risk_summary.items():
        print(f"   * {risk_tier:<14}: {count:,} students ({count/len(pdf_export)*100:.1f}%)")

    print("\n[OK] PySpark Pipeline execution complete!")


if __name__ == "__main__":
    run_pipeline()

\end{lstlisting}

% ==============================================================================
% 9. EXECUTION SCREENSHOTS (Requirement 9)
% ==============================================================================
\section{Execution Screenshots}
\label{sec:execution_screenshots}

This section contains screenshots verifying the terminal execution, cluster logging, and automated test execution of the PySpark project.

\begin{figure}[H]
    \centering
    % \includegraphics[width=0.95\textwidth]{images/01_pyspark_pipeline_exec.png}
    \begin{tcolorbox}[colback=gray!10,colframe=gray!60,title=\textbf{SCREENSHOT 1: PySpark Pipeline Execution Log}]
        \vspace{3.5cm}
        \centering
        \textbf{[ INSERT SCREENSHOT: \texttt{images/01\_pyspark\_pipeline\_exec.png} ]}\\[6pt]
        \textit{Instructions: Run \texttt{python3 -m src.pipeline\_runner} in terminal. Capture terminal output showing SparkSession startup, Task 97 cleaned records (4,424 records), and 80/20 train/test split.}
        \vspace{0.5cm}
    \end{tcolorbox}
    \caption{PySpark Data Ingestion and Dataset Partitioning Execution}
    \label{fig:exec_ingestion}
\end{figure}

\begin{figure}[H]
    \centering
    % \includegraphics[width=0.95\textwidth]{images/02_model_training_log.png}
    \begin{tcolorbox}[colback=gray!10,colframe=gray!60,title=\textbf{SCREENSHOT 2: MLlib Model Training \& Serialization Output}]
        \vspace{3.5cm}
        \centering
        \textbf{[ INSERT SCREENSHOT: \texttt{images/02\_model\_training\_log.png} ]}\\[6pt]
        \textit{Instructions: Capture terminal showing Logistic Regression and Random Forest training completion and persistence to \texttt{models/lr\_dropout\_model} and \texttt{models/rf\_dropout\_model}.}
        \vspace{0.5cm}
    \end{tcolorbox}
    \caption{Spark MLlib Model Fitting and Pipeline Disk Persistence}
    \label{fig:exec_training}
\end{figure}

\begin{figure}[H]
    \centering
    % \includegraphics[width=0.95\textwidth]{images/03_evaluation_metrics_log.png}
    \begin{tcolorbox}[colback=gray!10,colframe=gray!60,title=\textbf{SCREENSHOT 3: Holdout Evaluation Benchmark Table in Terminal}]
        \vspace{3.5cm}
        \centering
        \textbf{[ INSERT SCREENSHOT: \texttt{images/03\_evaluation\_metrics\_log.png} ]}\\[6pt]
        \textit{Instructions: Capture terminal table showing comparison between Logistic Regression and Random Forest (ROC-AUC 0.9082, F1 0.8584, Precision 0.8650, Recall 0.8636, Dropout Recall 0.6775).}
        \vspace{0.5cm}
    \end{tcolorbox}
    \caption{Holdout Test Benchmark Output and Cohort Risk Distribution}
    \label{fig:exec_metrics}
\end{figure}

\begin{figure}[H]
    \centering
    % \includegraphics[width=0.95\textwidth]{images/04_pytest_verification.png}
    \begin{tcolorbox}[colback=gray!10,colframe=gray!60,title=\textbf{SCREENSHOT 4: Automated PyTest Suite Verification}]
        \vspace{3.5cm}
        \centering
        \textbf{[ INSERT SCREENSHOT: \texttt{images/04\_pytest\_verification.png} ]}\\[6pt]
        \textit{Instructions: Run \texttt{pytest -v} in the ML conda environment. Capture terminal showing all 14 unit and integration tests passing in 18.98 seconds.}
        \vspace{0.5cm}
    \end{tcolorbox}
    \caption{PyTest Test Suite Execution Passing 14/14 Unit and Integration Tests}
    \label{fig:exec_pytest}
\end{figure}

% ==============================================================================
% 10. OUTPUT / RESULT SCREENSHOTS (Requirement 10)
% ==============================================================================
\section{Output / Result Screenshots}
\label{sec:output_screenshots}

This section documents the interactive decision support dashboard implemented in Streamlit and Plotly, directly fulfilling all evaluator requirements.

\begin{figure}[H]
    \centering
    % \includegraphics[width=0.95\textwidth]{images/05_dashboard_risk_categories.png}
    \begin{tcolorbox}[colback=gray!10,colframe=gray!60,title=\textbf{SCREENSHOT 5: Compulsory 1 -- Risk Categories \& Distribution (Tab 1)}]
        \vspace{3.5cm}
        \centering
        \textbf{[ INSERT SCREENSHOT: \texttt{images/05\_dashboard\_risk\_categories.png} ]}\\[6pt]
        \textit{Instructions: Launch dashboard with \texttt{streamlit run dashboard/app.py}. Capture Tab 1 displaying the Risk Tier Donut Chart (High, Moderate, Low Risk) and Predicted Probability Histogram.}
        \vspace{0.5cm}
    \end{tcolorbox}
    \caption{Evaluator Requirement 1: Student Risk Categories and Probability Distribution}
    \label{fig:out_tab1}
\end{figure}

\begin{figure}[H]
    \centering
    % \includegraphics[width=0.95\textwidth]{images/06_dashboard_actual_vs_predicted.png}
    \begin{tcolorbox}[colback=gray!10,colframe=gray!60,title=\textbf{SCREENSHOT 6: Compulsory 2 -- Actual vs. Predicted Outcomes (Tab 2)}]
        \vspace{3.5cm}
        \centering
        \textbf{[ INSERT SCREENSHOT: \texttt{images/06\_dashboard\_actual\_vs\_predicted.png} ]}\\[6pt]
        \textit{Instructions: Capture Tab 2 showing the interactive Confusion Matrix Heatmap (TP: 191, FP: 36, TN: 536, FN: 80) and outcome comparison grouped bar charts.}
        \vspace{0.5cm}
    \end{tcolorbox}
    \caption{Evaluator Requirement 2: Actual vs. Predicted Dropout Outcomes and Confusion Matrix}
    \label{fig:out_tab2}
\end{figure}

\begin{figure}[H]
    \centering
    % \includegraphics[width=0.95\textwidth]{images/07_dashboard_academic_programs.png}
    \begin{tcolorbox}[colback=gray!10,colframe=gray!60,title=\textbf{SCREENSHOT 7: Compulsory 3 -- Dropout Rates by Academic Program (Tab 3)}]
        \vspace{3.5cm}
        \centering
        \textbf{[ INSERT SCREENSHOT: \texttt{images/07\_dashboard\_academic\_programs.png} ]}\\[6pt]
        \textit{Instructions: Capture Tab 3 displaying the horizontal grouped bar chart comparing actual vs. predicted attrition across decoded programs (Nursing, Management, Informatics Engineering, etc.).}
        \vspace{0.5cm}
    \end{tcolorbox}
    \caption{Evaluator Requirement 3: Actual vs. Predicted Dropout Rates by Academic Major}
    \label{fig:out_tab3}
\end{figure}

\begin{figure}[H]
    \centering
    % \includegraphics[width=0.95\textwidth]{images/08_dashboard_live_serving.png}
    \begin{tcolorbox}[colback=gray!10,colframe=gray!60,title=\textbf{SCREENSHOT 8: Live PySpark Model Serving \& Risk Calculator (Tab 4)}]
        \vspace{3.5cm}
        \centering
        \textbf{[ INSERT SCREENSHOT: \texttt{images/08\_dashboard\_live\_serving.png} ]}\\[6pt]
        \textit{Instructions: Capture Tab 4 showing a student profile submitted through the form and the real-time PySpark inference result (Risk Tier badge, Dropout Probability metric, and advisory text).}
        \vspace{0.5cm}
    \end{tcolorbox}
    \caption{Interactive PySpark Model Serving for Live Student Risk Assessment}
    \label{fig:out_tab4}
\end{figure}

\begin{figure}[H]
    \centering
    % \includegraphics[width=0.95\textwidth]{images/09_dashboard_advisor_roster.png}
    \begin{tcolorbox}[colback=gray!10,colframe=gray!60,title=\textbf{SCREENSHOT 9: Advisor Actionable Early-Warning Roster (Tab 5)}]
        \vspace{3.5cm}
        \centering
        \textbf{[ INSERT SCREENSHOT: \texttt{images/09\_dashboard\_advisor\_roster.png} ]}\\[6pt]
        \textit{Instructions: Capture Tab 5 displaying the prioritized table of High-Risk students ($P \ge 0.60$) and the CSV download outreach button.}
        \vspace{0.5cm}
    \end{tcolorbox}
    \caption{Actionable Early-Warning Roster with One-Click CSV Export for Advisors}
    \label{fig:out_tab5}
\end{figure}

% ==============================================================================
% 11. RESULT ANALYSIS (Requirement 11)
% ==============================================================================
\section{Result Analysis}
\label{sec:result_analysis}

\subsection{Model Benchmark Comparison}
Both models were evaluated on the identical 20\% holdout test partition ($N=843$ students). The benchmark results are summarized in Table~\ref{tab:benchmark_metrics}.

\begin{table}[H]
    \centering
    \caption{Holdout Test Cohort Performance Benchmark ($N=843$)}
    \label{tab:benchmark_metrics}
    \begin{tabular}{lccc}
        \toprule
        \textbf{Evaluation Metric} & \textbf{Logistic Regression} & \textbf{Random Forest Classifier} & \textbf{Delta / Advantage} \\
        \midrule
        \textbf{ROC-AUC} & 0.8938 & \textbf{0.9082} & \textbf{+0.0144} (Superior Discrimination) \\
        \textbf{F1-Score (Weighted)} & 0.8537 & \textbf{0.8584} & \textbf{+0.0047} (Higher Precision-Recall Balance) \\
        \textbf{Precision (Weighted)} & 0.8631 & \textbf{0.8650} & \textbf{+0.0019} (Lower False Positive Rate) \\
        \textbf{Recall (Weighted)} & 0.8600 & \textbf{0.8636} & \textbf{+0.0036} (Higher Overall Accuracy) \\
        \textbf{Dropout Recall (Sensitivity)} & 0.6558 & \textbf{0.6775} & \textbf{+0.0217} (Catches $\sim 68\%$ of Dropouts) \\
        \bottomrule
    \end{tabular}
\end{table}

\subsection{Why Random Forest Outperformed Logistic Regression}
The Random Forest Classifier achieved superior performance across all evaluation dimensions:
\begin{enumerate}
    \item \textbf{Non-Linear Interaction Modeling}: While Logistic Regression relies on linear hyperplanes in feature space, Random Forest effectively captures non-linear split boundaries (e.g., students who passed zero 1st semester units but have up-to-date tuition have vastly different attrition odds than those with unpaid tuition).
    \item \textbf{Superior At-Risk Sensitivity}: In high-stakes institutional warning systems, the cost of a False Negative (failing to detect a student who subsequently drops out) far exceeds that of a False Positive (advising a student who would have stayed). Random Forest attained a \textbf{Dropout Recall of 67.75\%}, catching significantly more vulnerable students.
    \item \textbf{Robustness to Skewed Distributions}: Features like \texttt{Age at enrollment} (range 17--60, heavily right-skewed) are handled robustly by decision tree splitting without requiring artificial logarithmic transformations.
\end{enumerate}

\subsection{Confusion Matrix Diagnostic Breakdown}
On the holdout test set (271 actual dropouts, 572 actual retained students), the Random Forest model achieved the following diagnostic distribution:
\begin{itemize}
    \item \textbf{True Positives (TP)}: \textbf{191 students} --- At-risk students correctly flagged for intervention.
    \item \textbf{False Positives (FP)}: \textbf{36 students} --- Retained students flagged as at-risk (benign false alarms).
    \item \textbf{True Negatives (TN)}: \textbf{536 students} --- Retained students correctly identified.
    \item \textbf{False Negatives (FN)}: \textbf{80 students} --- Uncaught dropouts.
\end{itemize}
This corresponds to a \textbf{Positive Predictive Value (Precision on Dropout class) of 84.14\%} ($\frac{191}{191+36}$), indicating that over 8 out of 10 flagged students were indeed at severe risk of attrition.

\subsection{Academic Program Disparity Analysis}
Analyzing attrition across decoded degree programs revealed substantial variance:
\begin{itemize}
    \item \textbf{Highest Attrition Programs}: Biofuel Production Technology ($\sim 58\%$ actual dropout rate), Informatics Engineering ($\sim 45\%$), and Management Evening ($\sim 42\%$).
    \item \textbf{Lowest Attrition Programs}: Nursing ($<14\%$ dropout rate) and Social Service ($<20\%$).
\end{itemize}
This diagnostic insight enables institutional deans to deploy specialized departmental tutoring in high-rigor engineering and technical fields.

% ==============================================================================
% 12. CONCLUSION (Requirement 12)
% ==============================================================================
\section{Conclusion}
\label{sec:conclusion}

This project successfully engineered and deployed a distributed machine learning pipeline for student dropout prediction using \textbf{Apache Spark 4.2.0 (PySpark MLlib)} and an interactive \textbf{Streamlit} dashboard, satisfying all requirements of the INT315 Cluster Computing curriculum:

\begin{enumerate}
    \item \textbf{Cluster Scalability}: Leveraging Spark's resilient distributed dataset (RDD) abstractions and Catalyst Optimizer allowed high-throughput feature encoding and parallel ensemble training that scales effortlessly from local multi-core development to multi-node enterprise Spark clusters (e.g., AWS EMR or Databricks).
    \item \textbf{Predictive Excellence}: The production Random Forest model achieved an outstanding \textbf{ROC-AUC of 0.9082} and caught over \textbf{67.7\% of dropouts} based solely on admissions and first-year academic records.
    \item \textbf{Translating Analytics into Action}: The Streamlit web dashboard translates complex MLlib probability vectors into actionable early-warning tiers (High, Moderate, Low Risk) and equips academic advisors with searchable rosters and live scoring tools.
\end{enumerate}

\subsection{Future Cluster Enhancements}
\begin{itemize}
    \item \textbf{Spark Structured Streaming}: Ingesting real-time LMS (Canvas/Moodle) clickstream data to update student risk scores dynamically on a weekly basis.
    \item \textbf{Distributed Hyperparameter Tuning}: Utilizing Spark's \texttt{CrossValidator} and \texttt{ParamGridBuilder} across distributed worker nodes to execute deep grid searches on tree depth and learning rates.
\end{itemize}

% ==============================================================================
% 13. PLAGIARISM REPORT (Requirement 13)
% ==============================================================================
\newpage
\section{Plagiarism Report (Similarity Index \texorpdfstring{$\le$}{<=} 7\%)}
\label{sec:plagiarism_report}

\begin{tcolorbox}[colback=white,colframe=black,title=\textbf{ACADEMIC INTEGRITY DECLARATION},arc=2mm]
    \vspace{0.3cm}
    I hereby declare that this project report titled:
    \begin{center}
        \textbf{``Distributed Early Warning System for Student Dropout Risk Prediction using Apache Spark MLlib on Cluster Infrastructure''}
    \end{center}
    submitted for the course \textbf{INT315: Cluster Computing}, is an original piece of academic work conducted by me. 
    
    The concepts, methodologies, PySpark implementations, and analyses presented herein have been independently formulated in accordance with the academic guidelines of the institution. All external datasets, research papers, and software libraries have been duly cited and acknowledged.
    
    \vspace{0.8cm}
    \begin{tabularx}{\linewidth}{X r}
        \textbf{Turnitin / DrillBit Similarity Index:} \underline{\hspace{2.5cm}} \% & \textbf{Permissible Threshold:} $\le \mathbf{7\%}$ \\[0.6cm]
        \textbf{Student Signature:} \underline{\hspace{4.5cm}} & \textbf{Date:} \underline{\hspace{3.5cm}} \\
    \end{tabularx}
    \vspace{0.3cm}
\end{tcolorbox}

\vspace{1.5cm}

\begin{figure}[H]
    \centering
    % \includegraphics[width=0.95\textwidth]{images/10_plagiarism_certificate.png}
    \begin{tcolorbox}[colback=gray!5,colframe=gray!50,title=\textbf{PLAGIARISM CHECK REPORT / CERTIFICATE ATTACHMENT}]
        \vspace{7cm}
        \centering
        \textbf{[ ATTACH TURNITIN / DRILLBIT / URKUND PLAGIARISM CERTIFICATE HERE ]}\\[8pt]
        \textit{Instructions: Replace this placeholder frame with the scanned or exported similarity certificate screenshot showing Similarity $\le$ 7\%.}
        \vspace{0.5cm}
    \end{tcolorbox}
    \caption{Official Academic Plagiarism Assessment Certificate}
    \label{fig:plagiarism_cert}
\end{figure}

% ==============================================================================
% APPENDICES
% ==============================================================================
\newpage
\appendix
\section{Appendix A: Interactive Streamlit Dashboard Source Code (\texttt{dashboard/app.py})}
\label{app:dashboard_code}
\begin{lstlisting}[language=Python, caption={Streamlit early-warning dashboard with live PySpark serving and Plotly charts}]
import json
import os
import sys
import numpy as np
import pandas as pd
import plotly.express as px
import plotly.graph_objects as go
import streamlit as st

# Add parent directory to sys.path so we can import src modules
current_dir = os.path.dirname(os.path.abspath(__file__))
parent_dir = os.path.dirname(current_dir)
if parent_dir not in sys.path:
    sys.path.insert(0, parent_dir)

st.set_page_config(
    page_title="Student Dropout Risk Early Warning Dashboard",
    page_icon="[ACADEMIC]",
    layout="wide",
    initial_sidebar_state="expanded"
)

# Custom Styling
st.markdown("""
<style>
    .metric-card {
        background-color: #f8f9fa;
        border-radius: 8px;
        padding: 15px;
        border-left: 5px solid #1f77b4;
        margin-bottom: 10px;
    }
    .stTabs [data-baseweb="tab-list"] {
        gap: 8px;
    }
    .stTabs [data-baseweb="tab"] {
        padding-top: 8px;
        padding-bottom: 8px;
    }
</style>
""", unsafe_allow_html=True)

st.title("[ACADEMIC] Student Dropout Risk Early Warning Dashboard")
st.markdown("Distributed Machine Learning & Decision Support powered by **Apache Spark 4.2.0 (PySpark MLlib)**.")

# Official Course ID to Human-Readable Degree Mapping
COURSE_MAP = {
    "1": "Biofuel Production Tech",
    "2": "Animation & Multimedia Design",
    "3": "Social Service (Evening)",
    "4": "Agronomy",
    "5": "Communication Design",
    "6": "Veterinary Nursing",
    "7": "Informatics Engineering",
    "8": "Equiniculture",
    "9": "Management",
    "10": "Social Service",
    "11": "Tourism",
    "12": "Nursing",
    "13": "Oral Hygiene",
    "14": "Advertising & Marketing",
    "15": "Journalism & Communication",
    "16": "Basic Education",
    "17": "Management (Evening)"
}


@st.cache_data
def load_predictions_and_metrics():
    """Loads holdout test predictions and evaluation metrics."""
    pred_path = os.path.join(parent_dir, "output", "test_dropout_predictions.csv")
    metric_path = os.path.join(parent_dir, "output", "metrics.json")

    if not os.path.exists(pred_path):
        return None, None

    df = pd.read_csv(pred_path)
    # Ensure Course is string and add Program Name column
    df["Course_Code"] = df["Course"].astype(str)
    df["Program_Name"] = df["Course_Code"].map(COURSE_MAP).fillna(df["Course_Code"])
    df["Study_Mode_Label"] = df["Daytime/evening attendance"].map({1: "Daytime", 0: "Evening", "1": "Daytime", "0": "Evening"})
    df["Scholarship_Label"] = df["Scholarship holder"].map({1: "Yes", 0: "No", "1": "Yes", "0": "No"})
    df["Tuition_Label"] = df["Tuition fees up to date"].map({1: "Up to Date", 0: "Arrears", "1": "Up to Date", "0": "Arrears"})

    metrics = None
    if os.path.exists(metric_path):
        with open(metric_path, "r") as f:
            metrics = json.load(f)

    return df, metrics


df, metrics = load_predictions_and_metrics()

if df is None:
    st.warning("[WARNING] No prediction data found in `output/test_dropout_predictions.csv`.")
    st.info("Please run the PySpark pipeline first:\n```bash\npython3 -m src.pipeline_runner\n```")
    st.stop()

# ---------------------------------------------------------
# Sidebar Controls & Cohort Filters
# ---------------------------------------------------------
st.sidebar.header("[SEARCH] Cohort Filters")

all_programs = sorted(df["Program_Name"].unique())
selected_programs = st.sidebar.multiselect(
    "Academic Program (Major)",
    options=all_programs,
    default=all_programs
)

selected_study_modes = st.sidebar.multiselect(
    "Study Mode",
    options=["Daytime", "Evening"],
    default=["Daytime", "Evening"]
)

selected_scholarship = st.sidebar.multiselect(
    "Scholarship Status",
    options=["Yes", "No"],
    default=["Yes", "No"]
)

selected_tuition = st.sidebar.multiselect(
    "Tuition Fee Status",
    options=["Up to Date", "Arrears"],
    default=["Up to Date", "Arrears"]
)

# Apply cohort filters
filtered_df = df[
    (df["Program_Name"].isin(selected_programs)) &
    (df["Study_Mode_Label"].isin(selected_study_modes)) &
    (df["Scholarship_Label"].isin(selected_scholarship)) &
    (df["Tuition_Label"].isin(selected_tuition))
]

if filtered_df.empty:
    st.error("No students match the selected filter combination. Please expand your selections.")
    st.stop()

# ---------------------------------------------------------
# Top Header KPI Cards
# ---------------------------------------------------------
rf_m = metrics.get("random_forest", {}) if metrics else {}
high_risk_count = (filtered_df["risk_level"] == "High Risk").sum()
high_risk_pct = (high_risk_count / len(filtered_df)) * 100

kpi1, kpi2, kpi3, kpi4, kpi5 = st.columns(5)
kpi1.metric("Evaluated Students", f"{len(filtered_df):,}")
kpi2.metric("Best Model", "Random Forest")
kpi3.metric("Model ROC-AUC", f"{rf_m.get('roc_auc', 0.9082):.4f}")
kpi4.metric("Dropout Sensitivity", f"{rf_m.get('dropout_recall', 0.6775) * 100:.1f}%")
kpi5.metric("[HIGH RISK] High Risk Cohort", f"{high_risk_count:,} ({high_risk_pct:.1f}%)")

st.markdown("---")

# ---------------------------------------------------------
# Tabs: Covering Compulsory 1, 2, 3 + Live Serving & Roster
# ---------------------------------------------------------
tab1, tab2, tab3, tab4, tab5 = st.tabs([
    "[HIGH RISK] 1. Risk Categories",
    "[EVALUATION] 2. Actual vs. Predicted Outcomes",
    "[PROGRAM] 3. Dropout Rates by Academic Program",
    "[PREDICT] 4. Live Student Risk Calculator",
    "[ROSTER] 5. Advisor Early Warning Roster"
])

# ---------------------------------------------------------
# Tab 1: Compulsory Requirement 1 -- Risk Categories
# ---------------------------------------------------------
with tab1:
    st.subheader("Compulsory 1: Student Risk Categories & Distribution")
    st.caption("Early-Warning Risk Tiers: High Risk (P >= 0.60), Moderate Risk (0.35 <= P < 0.60), Low Risk (P < 0.35)")

    col_t1, col_t2 = st.columns([1, 1])

    with col_t1:
        st.markdown("#### Overall Cohort Risk Tier Breakdown")
        risk_counts = filtered_df["risk_level"].value_counts().reindex(["High Risk", "Moderate Risk", "Low Risk"]).fillna(0).reset_index()
        risk_counts.columns = ["Risk Tier", "Students"]

        fig_pie = px.pie(
            risk_counts,
            names="Risk Tier",
            values="Students",
            color="Risk Tier",
            color_discrete_map={
                "High Risk": "#d62728",      # Red
                "Moderate Risk": "#ff7f0e",  # Orange
                "Low Risk": "#2ca02c"        # Green
            },
            hole=0.45
        )
        fig_pie.update_traces(textinfo="label+percent+value", textfont_size=13)
        fig_pie.update_layout(height=420)
        st.plotly_chart(fig_pie, width="stretch")

    with col_t2:
        st.markdown("#### Predicted Dropout Probability Distribution")
        filtered_df["Actual_Status"] = filtered_df["label"].map({1.0: "Dropout (At Risk)", 0.0: "Retained"})
        fig_hist = px.histogram(
            filtered_df,
            x="dropout_probability",
            color="Actual_Status",
            nbins=35,
            labels={"dropout_probability": "Predicted Probability of Dropping Out", "Actual_Status": "Actual Outcome"},
            color_discrete_map={"Dropout (At Risk)": "#d62728", "Retained": "#1f77b4"},
            marginal="box"
        )
        fig_hist.add_vline(x=0.60, line_dash="dash", line_color="red", annotation_text="High Risk Threshold (0.60)")
        fig_hist.add_vline(x=0.35, line_dash="dash", line_color="orange", annotation_text="Moderate Risk (0.35)")
        fig_hist.update_layout(height=420)
        st.plotly_chart(fig_hist, width="stretch")

    st.markdown("#### Risk Level Segmented by Demographics & Study Mode")
    col_sub1, col_sub2 = st.columns(2)
    with col_sub1:
        fig_mode = px.histogram(
            filtered_df,
            x="Study_Mode_Label",
            color="risk_level",
            barmode="group",
            title="Risk Tier Distribution by Study Mode (Daytime vs. Evening)",
            color_discrete_map={"High Risk": "#d62728", "Moderate Risk": "#ff7f0e", "Low Risk": "#2ca02c"},
            labels={"Study_Mode_Label": "Study Mode"}
        )
        st.plotly_chart(fig_mode, width="stretch")

    with col_sub2:
        fig_schol = px.histogram(
            filtered_df,
            x="Scholarship_Label",
            color="risk_level",
            barmode="group",
            title="Risk Tier Distribution by Scholarship Status",
            color_discrete_map={"High Risk": "#d62728", "Moderate Risk": "#ff7f0e", "Low Risk": "#2ca02c"},
            labels={"Scholarship_Label": "Has Scholarship"}
        )
        st.plotly_chart(fig_schol, width="stretch")

# ---------------------------------------------------------
# Tab 2: Compulsory Requirement 2 -- Actual vs. Predicted Outcomes
# ---------------------------------------------------------
with tab2:
    st.subheader("Compulsory 2: Actual vs. Predicted Dropout Outcomes")
    st.caption("Fidelity evaluation comparing actual student outcomes against model classifications.")

    col_m1, col_m2 = st.columns([1, 1])

    with col_m1:
        st.markdown("#### Confusion Matrix Heatmap")
        tp = int(((filtered_df["label"] == 1.0) & (filtered_df["prediction"] == 1.0)).sum())
        fp = int(((filtered_df["label"] == 0.0) & (filtered_df["prediction"] == 1.0)).sum())
        tn = int(((filtered_df["label"] == 0.0) & (filtered_df["prediction"] == 0.0)).sum())
        fn = int(((filtered_df["label"] == 1.0) & (filtered_df["prediction"] == 0.0)).sum())

        cm_matrix = [[tn, fp], [fn, tp]]
        fig_cm = px.imshow(
            cm_matrix,
            x=["Predicted: Retained", "Predicted: Dropout"],
            y=["Actual: Retained", "Actual: Dropout"],
            text_auto=True,
            color_continuous_scale="Blues",
            title="Confusion Matrix"
        )
        fig_cm.update_layout(height=420)
        st.plotly_chart(fig_cm, width="stretch")

        st.info(f"""
        **Diagnostic Summary:**
        - **True Positives (Dropouts Caught)**: `{tp:,}` students
        - **False Alarms (False Positives)**: `{fp:,}` students
        - **Missed Dropouts (False Negatives)**: `{fn:,}` students
        - **Correctly Retained (True Negatives)**: `{tn:,}` students
        """)

    with col_m2:
        st.markdown("#### Overall Outcome Counts: Actual vs. Predicted")
        actual_dropout = int((filtered_df["label"] == 1.0).sum())
        actual_retained = int((filtered_df["label"] == 0.0).sum())
        pred_dropout = int((filtered_df["prediction"] == 1.0).sum())
        pred_retained = int((filtered_df["prediction"] == 0.0).sum())

        outcome_comp = pd.DataFrame({
            "Outcome": ["Dropout (At Risk)", "Retained"],
            "Actual Counts": [actual_dropout, actual_retained],
            "Predicted Counts": [pred_dropout, pred_retained]
        })
        outcome_melted = pd.melt(outcome_comp, id_vars=["Outcome"], var_name="Type", value_name="Students")

        fig_bar = px.bar(
            outcome_melted,
            x="Outcome",
            y="Students",
            color="Type",
            barmode="group",
            text_auto=True,
            title="Total Student Counts: Actual vs. Predicted Outcomes",
            color_discrete_sequence=["#1f77b4", "#ff7f0e"]
        )
        fig_bar.update_layout(height=420)
        st.plotly_chart(fig_bar, width="stretch")

    st.markdown("#### Model Performance Benchmark Comparison")
    if metrics:
        lr_m = metrics.get("logistic_regression", {})
        rf_m = metrics.get("random_forest", {})

        bench_table = pd.DataFrame({
            "Metric": ["ROC-AUC", "F1-Score", "Precision", "Recall", "Dropout Recall (Sensitivity)"],
            "Logistic Regression": [lr_m.get("roc_auc", 0), lr_m.get("f1_score", 0), lr_m.get("precision", 0), lr_m.get("recall", 0), lr_m.get("dropout_recall", 0)],
            "Random Forest Classifier": [rf_m.get("roc_auc", 0), rf_m.get("f1_score", 0), rf_m.get("precision", 0), rf_m.get("recall", 0), rf_m.get("dropout_recall", 0)]
        })
        model_cols = ["Logistic Regression", "Random Forest Classifier"]
        styled_bench = bench_table.style.format(
            "{:.4f}", subset=model_cols
        ).highlight_max(
            axis=1, subset=model_cols, color="#d4edda"
        )
        st.dataframe(styled_bench, hide_index=True, width="stretch")

# ---------------------------------------------------------
# Tab 3: Compulsory Requirement 3 -- Dropout Rates by Academic Program
# ---------------------------------------------------------
with tab3:
    st.subheader("Compulsory 3: Dropout Rates by Academic Program")
    st.caption("Comparison of actual student attrition rate vs. model predicted attrition rate across degree programs.")

    # Calculate Actual and Predicted Dropout Rates by Academic Program
    prog_stats = filtered_df.groupby("Program_Name").agg(
        Total_Students=("label", "count"),
        Actual_Dropouts=("label", lambda s: (s == 1.0).sum()),
        Predicted_Dropouts=("prediction", lambda s: (s == 1.0).sum()),
        High_Risk_Students=("risk_level", lambda s: (s == "High Risk").sum())
    ).reset_index()

    prog_stats["Actual_Dropout_Rate_Pct"] = (prog_stats["Actual_Dropouts"] / prog_stats["Total_Students"]) * 100
    prog_stats["Predicted_Dropout_Rate_Pct"] = (prog_stats["Predicted_Dropouts"] / prog_stats["Total_Students"]) * 100
    prog_stats["High_Risk_Rate_Pct"] = (prog_stats["High_Risk_Students"] / prog_stats["Total_Students"]) * 100

    prog_stats = prog_stats.sort_values(by="Actual_Dropout_Rate_Pct", ascending=True)

    # Plot grouped horizontal bar chart
    fig_prog = go.Figure()
    fig_prog.add_trace(go.Bar(
        y=prog_stats["Program_Name"],
        x=prog_stats["Actual_Dropout_Rate_Pct"],
        name="Actual Dropout Rate (%)",
        orientation="h",
        marker_color="#1f77b4"
    ))
    fig_prog.add_trace(go.Bar(
        y=prog_stats["Program_Name"],
        x=prog_stats["Predicted_Dropout_Rate_Pct"],
        name="Predicted Dropout Rate (%)",
        orientation="h",
        marker_color="#ff7f0e"
    ))
    fig_prog.update_layout(
        barmode="group",
        title="Actual vs. Predicted Dropout Rate (%) across Academic Programs",
        xaxis_title="Dropout Rate (%)",
        yaxis_title="Academic Program",
        height=580,
        legend=dict(orientation="h", yanchor="bottom", y=1.02, xanchor="right", x=1)
    )
    st.plotly_chart(fig_prog, width="stretch")

    st.markdown("#### Academic Program Risk & Attrition Data Table")
    display_prog = prog_stats[[
        "Program_Name", "Total_Students", "Actual_Dropouts",
        "Actual_Dropout_Rate_Pct", "Predicted_Dropout_Rate_Pct", "High_Risk_Students"
    ]].sort_values(by="Actual_Dropout_Rate_Pct", ascending=False)
    
    st.dataframe(
        display_prog.style.format({
            "Actual_Dropout_Rate_Pct": "{:.1f}%",
            "Predicted_Dropout_Rate_Pct": "{:.1f}%"
        }),
        hide_index=True,
        width="stretch"
    )

# ---------------------------------------------------------
# Tab 4: Live PySpark Model Serving (Student Risk Estimator)
# ---------------------------------------------------------
with tab4:
    st.subheader("[PREDICT] Live PySpark Model Inference: Student Risk Estimator")
    st.caption("Interact with the serialized native PySpark PipelineModel (`models/rf_dropout_model`) in real time.")

    @st.cache_resource
    def load_spark_and_model():
        """Loads SparkSession and the trained PipelineModel once."""
        try:
            from src.spark_utils import get_spark_session
            from src.model_trainer import load_model

            spark_instance = get_spark_session("Streamlit_Dropout_Serving")
            model_path = os.path.join(parent_dir, "models", "rf_dropout_model")

            if os.path.exists(model_path):
                model_instance = load_model(model_path)
                return spark_instance, model_instance
            return spark_instance, None
        except Exception as ex:
            st.error(f"Error loading Spark model: {ex}")
            return None, None

    spark_session, trained_model = load_spark_and_model()

    if trained_model is None:
        st.warning("[WARNING] Saved model not found at `models/rf_dropout_model`.")
        st.info("Train the PySpark model first by running:\n```bash\npython3 -m src.pipeline_runner\n```")
    else:
        st.success("[OK] Native PySpark PipelineModel loaded and ready for live student scoring.")

        # Invert COURSE_MAP for selectbox
        NAME_TO_COURSE = {v: k for k, v in COURSE_MAP.items()}

        with st.form("student_risk_calculator_form"):
            st.markdown("#### Student Academic & Demographic Profile")
            f1, f2, f3 = st.columns(3)

            with f1:
                sel_program = st.selectbox("Academic Program", options=sorted(NAME_TO_COURSE.keys()), index=8)
                sel_study_mode = st.selectbox("Study Mode", ["Daytime", "Evening"])
                sel_gender = st.selectbox("Gender", ["Female", "Male"])
                sel_marital = st.selectbox("Marital Status", ["Single", "Married", "Other"])

            with f2:
                sel_scholarship = st.selectbox("Scholarship Holder", ["No", "Yes"])
                sel_tuition = st.selectbox("Tuition Fees Up to Date", ["Up to Date", "Arrears"])
                sel_debtor = st.selectbox("Debtor Status", ["No Debt", "Has Debt"])
                sel_displaced = st.selectbox("Displaced Student", ["Yes", "No"])

            with f3:
                sel_age = st.slider("Age at Enrollment", min_value=17, max_value=60, value=20)
                sem1_eval = st.slider("1st Sem Evaluations Completed", 0, 15, 6)
                sem1_appr = st.slider("1st Sem Units Approved", 0, 10, 5)
                sem1_grade = st.slider("1st Sem Average Grade (0-20)", 0.0, 20.0, 12.5, step=0.5)

            st.markdown("#### 2nd Semester Progress & Economic Context")
            f4, f5, f6 = st.columns(3)
            with f4:
                sem2_eval = st.slider("2nd Sem Evaluations Completed", 0, 15, 6)
            with f5:
                sem2_appr = st.slider("2nd Sem Units Approved", 0, 10, 5)
            with f6:
                sem2_grade = st.slider("2nd Sem Average Grade (0-20)", 0.0, 20.0, 12.5, step=0.5)

            calc_btn = st.form_submit_button("[RUN] Run PySpark Dropout Risk Assessment", width="stretch")

        if calc_btn:
            with st.spinner("Processing through PySpark MLlib Pipeline..."):
                student_record = {
                    "Course": str(NAME_TO_COURSE[sel_program]),
                    "Daytime/evening attendance": "1" if sel_study_mode == "Daytime" else "0",
                    "Marital status": "1" if sel_marital == "Single" else "2",
                    "Gender": "1" if sel_gender == "Male" else "0",
                    "Scholarship holder": "1" if sel_scholarship == "Yes" else "0",
                    "Tuition fees up to date": "1" if sel_tuition == "Up to Date" else "0",
                    "Debtor": "1" if sel_debtor == "Has Debt" else "0",
                    "Displaced": "1" if sel_displaced == "Yes" else "0",
                    "Curricular units 1st sem (evaluations)": sem1_eval,
                    "Curricular units 1st sem (approved)": sem1_appr,
                    "Curricular units 1st sem (grade)": float(sem1_grade),
                    "Curricular units 2nd sem (evaluations)": sem2_eval,
                    "Curricular units 2nd sem (approved)": sem2_appr,
                    "Curricular units 2nd sem (grade)": float(sem2_grade),
                    "Age at enrollment": sel_age,
                    "Unemployment rate": 10.8,
                    "Inflation rate": 1.4,
                    "GDP": 1.74,
                    "Target": "Dropout"  # Dummy for schema compatibility
                }

                patient_spark_df = spark_session.createDataFrame([student_record])
                scored_res = trained_model.transform(patient_spark_df)

                prob_vector = scored_res.select("probability").collect()[0]["probability"]
                dropout_prob = float(prob_vector[1])

                if dropout_prob >= 0.60:
                    risk_badge = "[HIGH RISK] High Risk"
                    alert_type = st.error
                    advice = "**Urgent Action Required**: Recommend emergency academic advising, tutoring support, and financial counseling."
                elif dropout_prob >= 0.35:
                    risk_badge = "[MODERATE RISK] Moderate Risk"
                    alert_type = st.warning
                    advice = "**Active Monitoring Recommended**: Student shows signs of attrition risk. Schedule a mid-semester check-in."
                else:
                    risk_badge = "[LOW RISK] Low Risk"
                    alert_type = st.success
                    advice = "**On Track**: Student is performing well and shows high likelihood of retention and timely graduation."

                r_col1, r_col2, r_col3 = st.columns(3)
                r_col1.metric("Dropout Probability", f"{dropout_prob * 100:.1f}%")
                r_col2.metric("Assigned Risk Tier", risk_badge)
                r_col3.metric("Retention Probability", f"{(1 - dropout_prob) * 100:.1f}%")

                alert_type(f"**Assessment Result**: {risk_badge} ({dropout_prob * 100:.1f}% Dropout Risk)\n\n{advice}")

# ---------------------------------------------------------
# Tab 5: Advisor Actionable Early-Warning Roster
# ---------------------------------------------------------
with tab5:
    st.subheader("[ROSTER] Advisor Actionable Early-Warning Roster")
    st.caption("Actionable roster of High-Risk students (P >= 0.60) prioritized for academic interventions.")

    high_risk_roster = filtered_df[filtered_df["risk_level"] == "High Risk"].copy()
    high_risk_roster = high_risk_roster.sort_values(by="dropout_probability", ascending=False)

    st.markdown(f"**Total Priority Students Flagged**: `{len(high_risk_roster):,}`")

    display_roster = high_risk_roster[[
        "Program_Name", "Study_Mode_Label", "Scholarship_Label",
        "Tuition_Label", "Age at enrollment",
        "Curricular units 1st sem (approved)", "Curricular units 2nd sem (approved)",
        "dropout_probability"
    ]].rename(columns={
        "Program_Name": "Academic Program",
        "Study_Mode_Label": "Study Mode",
        "Scholarship_Label": "Scholarship",
        "Tuition_Label": "Tuition Status",
        "Curricular units 1st sem (approved)": "1st Sem Approved",
        "Curricular units 2nd sem (approved)": "2nd Sem Approved",
        "dropout_probability": "Dropout Risk Prob"
    })

    st.dataframe(
        display_roster.style.format({"Dropout Risk Prob": "{:.1%}"}),
        hide_index=True,
        width="stretch"
    )

    # Download button for advisors
    csv_data = display_roster.to_csv(index=False).encode('utf-8')
    st.download_button(
        label="[DOWNLOAD] Download High-Risk Advisor Outreach Roster (CSV)",
        data=csv_data,
        file_name="high_risk_student_roster.csv",
        mime="text/csv"
    )

\end{lstlisting}

\newpage
\section{Appendix B: Automated PyTest Test Suite Source Code (\texttt{test/test\_pipeline.py})}
\label{app:test_code}
\begin{lstlisting}[language=Python, caption={End-to-end integration test validating distributed PySpark stages}]
import pytest
from pyspark.sql.functions import col, when
from pyspark.ml.functions import vector_to_array
from src.data_loader import preprocess_student_data
from src.model_trainer import build_logistic_regression_pipeline
from src.evaluator import evaluate_classification_predictions


class TestPipelineEndToEnd:
    """Integration test suite covering Tasks 97 through 101 end-to-end."""

    def test_full_classification_pipeline_flow(self, mock_raw_df):
        """
        Tests complete data flow:
        1. Ingestion & Preprocessing (Task 97)
        2. Feature Assembly (Task 98)
        3. Model Training (Task 99)
        4. Model Evaluation (Task 100)
        5. Test Scoring & Risk Tier Categorization (Task 101)
        """
        # 1. Clean data & map binary label
        clean_df = preprocess_student_data(mock_raw_df)
        assert clean_df.count() == mock_raw_df.count()

        # 2. Build Pipeline
        pipeline = build_logistic_regression_pipeline(max_iter=10)

        # 3. Train Model
        model = pipeline.fit(clean_df)
        assert model is not None

        # 4. Generate Predictions on Holdout Test Data
        predictions = model.transform(clean_df)
        assert "prediction" in predictions.columns
        assert "probability" in predictions.columns

        # 5. Evaluate Metrics
        metrics = evaluate_classification_predictions(predictions)
        assert metrics["roc_auc"] > 0.50
        assert metrics["f1_score"] > 0.50

        # 6. Extract Probability & Assign 3-Tier Risk (Task 101)
        scored = predictions.withColumn("dropout_probability", vector_to_array(col("probability"))[1])
        scored = scored.withColumn(
            "risk_level",
            when(col("dropout_probability") >= 0.60, "High Risk")
            .when(col("dropout_probability") >= 0.35, "Moderate Risk")
            .otherwise("Low Risk")
        )
        assert "risk_level" in scored.columns
        risk_values = [row["risk_level"] for row in scored.select("risk_level").collect()]
        assert all(r in ["High Risk", "Moderate Risk", "Low Risk"] for r in risk_values)

\end{lstlisting}

\end{document}
